\chapter{Language Reference}

\section{Types}

Zinc is statically typed: every expression has a type that is known at compile time.

\subsection{Primitive Types}

\begin{center}
\begin{tabular}{ll}
\toprule
Type & Description \\
\midrule

\zn{u0} & Empty type (equivalent to void)\\
\zn{u1} & Boolean values\\
\zn{u8}, \zn{u16}, \zn{u32}, \zn{u64} & Unsigned integer with specific bit-width\\
\zn{i8}, \zn{i16}, \zn{i32}, \zn{i64} & Signed integer with specific bit-width\\
\zn{f32}, \zn{f64} & Floating point numbers\\
\bottomrule

\end{tabular}

\end{center}

\subsection{Compound Types}

\begin{center}
\begin{tabular}{ll}

\toprule
Type & Description\\

\midrule
\texttt{*T} & Pointer to \texttt{T}\\
\texttt{[]T} & Array of \texttt{T}\\
\texttt{[n]T} & Array of \texttt{T} with fixed size (n)\\
\texttt{(T1, T2)} & Tuples\\
\texttt{R (A, B)} & Function that returns \texttt{R} and takes \texttt{A}, \texttt{B}\\
\bottomrule

\end{tabular}
\end{center}

\section{Literals}

\begin{center}
\begin{tabular}{ll}

\toprule
Literal type & Value\\

\midrule
Integers & Decimal: \texttt{123}, Hex: \texttt{0xDEADBEEF}, binary: \texttt{0b100010}\\
Floats & Standard \texttt{3.14},  Scientific: \texttt{1e-10}\\
Boolean & \zn{true}, \zn{false}\\
Null pointer & \zn{none}\\
Strings & \texttt{*char}, content between quotes.\\

\bottomrule

\end{tabular}
\end{center}

\subsection{String Escapes}

\begin{itemize}
    \item{n} new line
    \item{t} tab
    \item{r} cursor at the start of the line
    \item{", '} These are inserted directly
    \item{uXXXX} unicode code point of 4 hex values
    \item{UXXXXXXXX} unicode code point of 8 hex values
\end{itemize}

\subsection{Arrays and Slices}

The unique difference between arrays and slices is that arrays have a fixed size known at compile time meanhwile slices have the size assigned at runtime.
But arrays and slices share the same memory layout and both hold the length of their elements.

\begin{lstlisting}[style=zinc, caption={Arrays and slices}]
[5]i32 nums = [1, 2, 3, 4, 5] // fixed-size (5)

mid := nums[1:4] // from 1 to 3
head := nums[:3] // from 0 to 3
tail := nums[3:] // from 3 to end

i := 0
for i < nums.len {
    nums[i] = i
    i = i + 1
}
\end{lstlisting}

\subsection{Tuples}
Tuples group different values into a single unnamed struct.

\begin{lstlisting}[style=zinc, caption={Tuple construction and access}]
pair := (12, "hello")
triple := (12, "hello", 3.14)

u32 val := triple.0
\end{lstlisting}

Tuples are the idiomatic way to return multiple values:

\begin{lstlisting}[style=zinc, caption={Multiple return values}]
divmod :: (u32, u32) (u32 a, u32 b) {
    return (a / b, a \% b)
}
\end{lstlisting}

\subsection{Enums}
Enums are tagged unions (sum types) that can carry per-variant data.
Each variant is a struct prefixed by a \texttt{u8} discriminant tag.

\begin{lstlisting}[style=zinc, caption={Enum declaration and construction}]
Shape :: enum {
    Square(f32)
    Rect(f32, f32)
}
square := Shape::Square(10.0)
rect := Shape::Rect(10.0, 12.0)
\end{lstlisting}

\section{Variables}

Variables live on the stack and are always mutable.
The type can be stated explicitly or inferred from the right-hand size (RHS).

\begin{lstlisting}[style=zinc, caption={Variable declarations}]
i32 val = 0 // explicit type
other := 10 // inferred type
\end{lstlisting}

\textbf{Destructuring} binds struct fields or tuple elements directly to names:

\begin{lstlisting}[style=zinc, caption={Destructuring}]
{x, y}          := point
{x: px, y: py}  := point // same with renaming x to px and y to py

(first, second) := pair

// Both way combined
({x, y}, {x: otherX, y: otherY}) := pair
\end{lstlistng}

\section{Functions}

\subsection{Declarations}
A normal function is composed by:
\zninline{<name> :: Ret (Args...)[Capabilities] \{ ... \}}
For short functions instead of declaring a block with a single return statement there is the \zn{->} operator that returns directly.

\begin{lstlisting}[style=zinc, caption={Function declarations}]

add :: u32 (u32 a, u32 b) {
    return a + b
}

mul :: u32 (u32 a, u32 b) -> a * b

main :: i32 () {
    return add(10, mul(12, 3))
}
\end{lstlistng}

\begin{note}
This section doesn't show an exaple with capabilities argument list.
For more information about capabilities see the section \ref{sec:capabilities}
\end{note}

\subsection{Receiver Methods}
The \zn{for} block attaches methods to a type.
The named binding after the type is the implicit \emph{self} argument

\begin{lstlisting}[style=zinc, caption{Receiver methods}]
i32 self :: for {
    add :: i32 (i32 other) -> return self + other

    abs :: u32 () {
        if self > 0 { return self as u32 }
        return (-self) as u32
    }
}

*i32 ptr :: for {
    incr :: u0 {
        *ptr = *ptr + 1
    }
}

main :: i32 () {
    val := 10
    val = val.add(5)
    val = 10.add(val)

    p := \&val
    p.incr()
    return 0
}
\end{lstlisting}

\subsection{Static Functions}

Omitting the receiver binding creates a static (namespace-scoped) function called with \zn{Type::name()} syntax.

\begin{lstlisting}[style=zinc, caption={Static functions}]
Counter :: struct {
    i32 count
}

Counter :: for {
    new :: Counter () -> Counter { count: 0 }
    from_val :: Counter (i32 val) -> Counter { count: val }
}

main :: i32 () {
    c := counter::new()

    c2 := Counter::from_val(10)
    return 0
}
\end{lstlisting}

Static functions are usually used as a constructors for used-defined types.

\section{Capabilities}
\label{sec:capabilities}

A capability is a contextual dependency declared in a second parameter list after the main one.
The called does not pass it explicitly, isntead it is threaded automatically from a surrounding \zn{with} block.

\begin{lstlisting}[style=zinc, caption={Capability declaration and use}]
use <core::mem>

allocBuf :: *u8 (u64 size)[Allocator alloc] {
    return alloc.alloc(size)
}

main :: i32 () {
    with _ := Allocator::heap() {
        buf := allocBuf(1024) // the allocator is automatically injected
    }

    allocBuf() // Compiler rejects this, Capability 'Allocator' doesn't exist
    return 0
}
\end{lstlistng}

Capabilities are propagated transitibely: a function that calls another capability-requiring function inherits the requirement automatically.

\begin{lstlisting}[style=zinc, caption={Transitive capability}]
allocBuf :: *u8 ()[Allocator allocator] {
    return allocator.alloc(1024)
}

transitive :: *u8 ()[Allocator] { // Unnamed capability
    return new() // Capability is passed implicitly
}

main :: i32 () {
    with _ := Allocator::heap() {
        transitive()
    }
    return 0
}
\end{lstlisting}

Transitive capabilities are often unused in the function itself but just passed to other functions (like in this example the \texttt{transitive} function doesn't use the \texttt{Allocator}) So the capability name can be omitted.

\subsection{Foreign Declarations}

Foreign declarations import C functions for use as ordinary Zinc functions.
Variadic C functions are supported with \texttt{...}.
Related functions can be grouped under a named namespace.

\begin{lstlisting}[style=zinc, caption={Foreign declarations}]
foreign printf :: u0 (*char, ...)

libc :: foreign {
    malloc :: *u8 (u64)
    free :: u0 (*u8)
}

main :: i32 () {
    printf("Hello world\n")
    buf := libc::malloc(1024)
    libc::free(buf)
}
\end{lstlisting}

There's also an example of a snake game using raylib in the repo.

\subsection{Control Flow}

\subsection{Conditionals}
The if/else statement can be used in three different ways:
- As a normal if statement
- As an expression using the \zn{->}
- To check the enum variant

\begin{lstlisting}[style=zinc, caption={if/else}]
// Normal if usage
if x > 0 {
    // positive
} else if x == 0 {
    // zero
} else {
    // negative
}

// Used as expression, the syntax is: if cond -> expr else -> expr
val := if x > 0 -> x else -> -x // used as abs

Expr :: enum {
    Add(*Expr, *Expr)
    Num(f64)
}

if Expr::Add(left, right) := myEnum {
    // use left and right here
} else if Option::Num(num) := myEnum {
    // use num here
}
\end{lstlisting}

\subsection{Loops}

Zinc uses \zn{for} for both C-style and while-style loops.
\begin{lstlisting}[style=zinc, caption={Loop forms}]
// While-style

i := 0
for i < 10 {
    i = i + 1
}

for j := 0; i < 10; j = j + 1 {
    // do something
}
\end{lstlisting}

\zn{break} exits the innermost loop and \zn{continue} skips to the next iteration.

\begin{lstlisting}[style=zinc, caption={break and continue}]
for i < 10 {
    i = i + 1
    if i == 5 { break }
    if i \% 2 == 0 { continue }
}
\end{lstlisting}

\subsection{Defer}

\zn{defer} schedules a statement to run when the enclosing scope exits, in last=in first-out order.

\begin{lstlisting}[style=zinc, caption={Defer}]
main :: i32 () {
    buf := libc::malloc(1024)
     defer libc::free(buf) // runs when the funciton returns

    return 0
}
\end{lstlisting}

\section{Operators}

\subsection{Update operator}
In zinc there isn't the classic \texttt{+=} or \texttt{-=} operators.
Instead iexists the operator \zn{:}. It reads a mutable location, transform its value through an expression and writes the result back.
So instead of having a lot of indipendent operators (\texttt{+=}, \texttt{-=}, texttt{|=} ...)
Zinc has only one operator that groups all of them.

\subsubsection*{Syntax}

\begin{center}
\begin{tabular}{ll}

\toprule
Form & Equivalent\\

\midrule
\zn{x : op expr} & \zn{x = x op expr}\\
\zn{x : unary\_op} & \zn{x = unary\_op x}\\
\zn{x : f(args)} & \zn{x = f(x, ...args)}\\
\zn{x : .method(args)} & \zn{x = x.method(args)}\\
\zn{x : .member} & \zn{x = x.member}\\

\bottomrule

\end{tabular}
\end{center}

The left-hand size (LHS) must be a mutable location: a variable, a field access,  a pointer reference or a subscript.
The right-hand side must be one of the four forms above. Plain assignment uses \zn{=} instead.

\subsubsection*{Basic usage}

\begin{lstlisting}[style=zinc, caption={Update operator}]
// Arithmetic - binary operators
count :+ 1
price :* 1.08
ratio :/ total

// Unary operator
flag : not
mask : ~

// Functions and method forms
buf     : compress()        // becomes compress(buf)
name    : .to_lower()       // name = name.to_lower()
value   : clamp(0, 100)     // value = clamp(value, 0, 100)
config  : merge(defaults)   // config = merge(config, defaults)
\end{lstlisting}

Since this operator is an assignment the RHS must produce the same type of the LHS.

\section{Type Aliases}
\begin{lstlisting}[style=zinc, caption={Type aliases}]
void :: type u0
bool :: type u1
int :: type i32
uint :: type u32

VeryComplicatedType :: (*i32(*char, *char), *MyStruct)
\end{lstlisting}

\section{Type Casting}

To cast a type zinc uses the \zn{as} keyword after the expression:
\begin{lstlisting}[style=zinc, caption={Type casting}]
foreign malloc :: *u8 (u64)

main :: i32 () {
    *MyStruct obj = malloc(sizeof MyStruct) // This is casted implicitly
    buff := malloc(sizeof MyStruct) as *MyStruct // explicitly casted from *u8 to *MyStruct
    return 0
}
\end{lstlisting}
The cast is possible only if the types are compatible.
For example a pointer cannot be casted to a number. An invalid cast produces a compiler error.

\section{Visitility}

In Zinc there are two level of visibility:
The default level means the symbol is visible everywhere in the same file.
By adding the \zn{pub} keyword that symbol becomes global and other files can use it.

\section{Imports}
Zinc uses the \zn{use} keyword for imports. They must be declared at file-level.

\begin{lstlisting}[style=zinc, caption={Imports}]
use myfile // imports a file in the same directory
use submodule::myfile // imports myfile under the submodule directory
use <core::mem> // <> means external libraries or std lib import
use super::sharedFile // super is a special keyword to go up the directory (translated to ..)
\end{lstlisting}


